\section{Future Work}
\label{sec:future}

The hostile audit concludes that P0 should proceed only after required protocol controls are frozen. The immediate next step is therefore corpus selection and preregistration under the audit gates, not execution.

P0 must first create a reproducible candidate task window, sequential-dependence rule, task hashes, canonical workspace materialisation procedure, FUTURE_HISTORY_LEAK_GATE, task-leakage audit, frozen hidden-verifier hashes, reconstruction-probe rubric, model/tool/environment manifest, stopping/rerun rules, and resource-classification rules.

After P0, repeated trials on the same complex repository are needed before any uncertainty estimate is meaningful. Public replications should then vary repository size, language, documentation quality, test quality, task locality, dependency width, and architecture.

Semantic-state ablation should remove small preregistered information classes rather than broadly degrading the repository. The scientific question is which durable artefacts have marginal continuity value.

Tiered execution belongs in a separate experiment because lower-cost exploration/execution is established prior work and changing executor capability would confound the core history treatment.

Cross-model continuity is particularly interesting but not part of P0. A design in which model family $M_A$ produces an accepted $R_t$ and a different family $M_B$ resumes from it without predecessor model-native state would provide stronger evidence that continuity resides in external project state rather than a model-specific session representation. It should be tested only after same-model causal identification works.

Reconstruction scaling should be tested directly. Synthetic irrelevant repository growth can hold task semantics constant; real tasks can vary dependency width. Plain search, symbol indexes, static-analysis graphs, repository maps, and model-guided retrieval can be compared without claiming any one is intrinsically RaS.

Security work should evaluate malicious repository instructions, tests, dependencies, commit metadata, poisoned durable documentation, work-package manipulation, executor compromise, and secret exposure. Least-privilege benefits remain architectural possibilities until adversarially measured.

The broader Authoritative-State Externalisation Principle remains speculative. Software repositories have unusually strong versioning, diffability, executable verification, and reproducible tooling; CRM systems, event logs, databases, and document stores may not offer equivalent semantics.
